% Biên dịch bằng XeLaTeX trên Overleaf.
\documentclass[12pt,a4paper,oneside]{report}
\usepackage[a4paper,top=2.2cm,bottom=2.2cm,left=3.2cm,right=2cm]{geometry}
\usepackage{fontspec}
\IfFontExistsTF{TeX Gyre Termes}{
  \setmainfont{TeX Gyre Termes}
  \setsansfont{TeX Gyre Heros}
  \setmonofont{DejaVu Sans Mono}[Scale=0.84]
}{
  \setmainfont{Times New Roman}
  \setsansfont{Arial}
  \setmonofont{Consolas}[Scale=0.84]
}
\usepackage[vietnamese]{babel}
\usepackage{microtype,setspace,graphicx,booktabs,longtable,tabularx,array,multirow}
\usepackage{enumitem,amsmath,amssymb,float,caption,xcolor,listings,tikz}
\usetikzlibrary{arrows.meta,positioning,shapes.geometric,fit}
\usepackage{fancyhdr,titlesec,tocloft,hyperref,url}
\onehalfspacing

\definecolor{schoolred}{HTML}{A61B1B}
\definecolor{schoolblue}{HTML}{173F7A}
\definecolor{codegreen}{HTML}{1B6E3C}
\definecolor{codegray}{HTML}{F4F5F7}
\hypersetup{colorlinks=true,linkcolor=schoolblue,citecolor=schoolblue,urlcolor=schoolblue,
  pdftitle={Báo cáo hệ thống thông tin chủ động cho hành động cá nhân},
  pdfauthor={Võ Trọng Tài và Nguyễn Cao Hậu}}
\pagestyle{fancy}
\fancyhf{}
\fancyhead[L]{\small Hệ thống thông tin chủ động cho hành động cá nhân}
\fancyhead[R]{\small\leftmark}
\fancyfoot[C]{\thepage}
\renewcommand{\headrulewidth}{0.4pt}
\setlength{\headheight}{15pt}
\titleformat{\chapter}[display]{\bfseries\centering}{\large CHƯƠNG \thechapter}{8pt}{\Large}
\titleformat{\section}{\bfseries\large}{\thesection}{0.7em}{}
\titleformat{\subsection}{\bfseries}{\thesubsection}{0.7em}{}
\setlist{nosep,leftmargin=1.7em}
\captionsetup{font=small,labelfont=bf}
\lstdefinestyle{source}{backgroundcolor=\color{codegray},basicstyle=\ttfamily\footnotesize,
 keywordstyle=\color{schoolblue}\bfseries,stringstyle=\color{codegreen},commentstyle=\color{gray},
 numbers=left,numberstyle=\tiny\color{gray},breaklines=true,frame=single,
 rulecolor=\color{gray!35},showstringspaces=false,tabsize=2,captionpos=b}
\lstset{style=source}

% Bổ sung MSSV còn thiếu tại đây trước khi nộp.
\newcommand{\studentOne}{Võ Trọng Tài}
\newcommand{\studentOneId}{2351220084}
\newcommand{\studentTwo}{Nguyễn Cao Hậu}
\newcommand{\studentTwoId}{\textit{Bổ sung MSSV}}
\newcommand{\advisor}{TS. Đỗ Phúc Hảo}
\newcommand{\placeholderfigure}[2]{\begin{figure}[H]\centering
 \IfFileExists{#1}{\includegraphics[width=0.78\textwidth,height=0.58\textheight,keepaspectratio]{#1}}
 {\fbox{\parbox[c][4.5cm][c]{0.76\textwidth}{\centering Bổ sung ảnh vào thư mục
 \texttt{images/}\\[4pt]\texttt{#1}}}}\caption{#2}\end{figure}}

\begin{document}
\pagenumbering{gobble}
\begin{titlepage}
\begin{tikzpicture}[remember picture,overlay]
 \draw[schoolred,line width=1.2pt] ([xshift=1.1cm,yshift=-1.1cm]current page.north west)
 rectangle ([xshift=-1.1cm,yshift=1.1cm]current page.south east);
 \draw[schoolred,line width=0.4pt] ([xshift=1.25cm,yshift=-1.25cm]current page.north west)
 rectangle ([xshift=-1.25cm,yshift=1.25cm]current page.south east);
\end{tikzpicture}
\begin{center}
{\bfseries BỘ GIÁO DỤC VÀ ĐÀO TẠO\\[4pt]}
{\bfseries\color{schoolred} TRƯỜNG ĐẠI HỌC KIẾN TRÚC ĐÀ NẴNG\\}
{\bfseries\color{schoolblue} KHOA CÔNG NGHỆ THÔNG TIN}
\vspace{1.3cm}
\IfFileExists{images/logo-truong.png}{\includegraphics[height=3cm]{images/logo-truong.png}}
 {\fbox{\parbox[c][2.5cm][c]{3.5cm}{\centering LOGO TRƯỜNG}}}
\vspace{1cm}
{\Large\bfseries BÁO CÁO ĐỒ ÁN MÔN HỌC\\[0.7cm]}
{\large\bfseries ĐỀ TÀI:\\[0.3cm]}
{\LARGE\bfseries\color{schoolred} XÂY DỰNG HỆ THỐNG THÔNG TIN CHỦ ĐỘNG\\
CHO HÀNH ĐỘNG CÁ NHÂN}
\vfill
\begin{tabular}{>{\bfseries}p{4.2cm}p{9.2cm}}
Sinh viên thực hiện: & \studentOne{} -- MSSV: \studentOneId\\
& \studentTwo{} -- MSSV: \studentTwoId\\[4pt]
Giảng viên hướng dẫn: & \advisor\\
Chuyên ngành: & Công nghệ thông tin\\
Công nghệ cốt lõi: & Flutter, Node.js/Express, MongoDB, AI và TTS tiếng Việt\\
Mã nguồn: & \url{https://github.com/Hau-altf4/action_planner}
\end{tabular}
\vfill
{\bfseries Đà Nẵng -- 2026}
\end{center}
\end{titlepage}

\clearpage\pagenumbering{roman}
\chapter*{LỜI CAM ĐOAN}\addcontentsline{toc}{chapter}{LỜI CAM ĐOAN}
Chúng tôi cam đoan báo cáo mô tả trung thực sản phẩm phần mềm do nhóm xây dựng.
Các công nghệ, thư viện mã nguồn mở và tài liệu tham khảo được ghi nhận tại cuối báo cáo.
Những số liệu kiểm thử chỉ sử dụng kết quả có thể tái lập từ mã nguồn và lệnh kiểm tra;
báo cáo không sử dụng các con số hiệu năng chưa được đo đạc.
\vspace{1.5cm}
\begin{flushright}Đà Nẵng, ngày \ldots{} tháng \ldots{} năm 2026\\
\textbf{Sinh viên thực hiện}\\[1.8cm]
\studentOne{} \hspace{1.4cm} \studentTwo\end{flushright}

\chapter*{LỜI CẢM ƠN}\addcontentsline{toc}{chapter}{LỜI CẢM ƠN}
Nhóm xin trân trọng cảm ơn \advisor{} đã hướng dẫn phương pháp phân tích hệ thống,
góp ý kiến trúc và hỗ trợ nhóm hoàn thiện đồ án. Nhóm cũng cảm ơn quý Thầy/Cô Khoa
Công nghệ Thông tin, Trường Đại học Kiến trúc Đà Nẵng đã trang bị kiến thức về lập
trình, cơ sở dữ liệu, trí tuệ nhân tạo và quy trình phát triển phần mềm.

\chapter*{TÓM TẮT ĐỒ ÁN}\addcontentsline{toc}{chapter}{TÓM TẮT ĐỒ ÁN}
Đồ án xây dựng hệ thống quản lý hành động cá nhân đa nền tảng, hỗ trợ tài khoản riêng,
quản lý vòng đời công việc, lập lịch theo khối thời gian, nhắc hạn, trợ lý AI, phản hồi,
quản trị và đọc nội dung bằng giọng nói tiếng Việt. Ứng dụng khách dùng Flutter; máy chủ
dùng Node.js, Express và MongoDB Native Driver; dữ liệu lưu trong MongoDB. Trợ lý AI
hỗ trợ Ollama, OpenAI, Google Gemini và Anthropic. Khóa API cá nhân được mã hóa
AES-256-GCM tại máy chủ. Tính năng giọng nói ưu tiên vietTTS và tự chuyển sang giọng
đọc hệ thống trên thiết bị khi mô hình cục bộ chưa sẵn sàng.

Hệ thống được kiểm tra bằng \texttt{flutter analyze}, bảy kiểm thử Flutter tự động,
kiểm tra cú pháp Node.js và quá trình tạo APK debug. Kết quả: không có lỗi phân tích
tĩnh, 7/7 kiểm thử đạt và APK được tạo thành công.

\textbf{Từ khóa:} Flutter, MongoDB, Node.js, quản lý công việc, AI, Ollama, vietTTS,
JWT, thông báo cục bộ.

\tableofcontents\listoffigures\listoftables
\clearpage\pagenumbering{arabic}

\chapter{MỞ ĐẦU}
\section{Đặt vấn đề}
Người dùng thường lưu công việc ở nhiều nơi như sổ tay, tin nhắn, email và ứng dụng
ghi chú. Dữ liệu phân tán làm tăng tải ghi nhớ, giảm khả năng xác định việc ưu tiên và
dễ dẫn đến bỏ lỡ thời hạn. Một hệ thống quản lý hành động cá nhân cần không chỉ ghi
nhận công việc mà còn hỗ trợ lập kế hoạch, nhắc hạn và đưa ra gợi ý có căn cứ.

\section{Mục tiêu}
Mục tiêu tổng quát là xây dựng ứng dụng đa nền tảng giúp người dùng quản lý công việc
từ lúc tạo, phân tích, lập lịch, theo dõi tiến độ đến hoàn thành, đồng thời bảo vệ dữ liệu
theo từng tài khoản. Các mục tiêu cụ thể gồm:
\begin{itemize}
 \item Xây dựng ứng dụng Flutter dùng chung mã nguồn cho nhiều nền tảng.
 \item Xây dựng REST API có JWT và phân quyền người dùng/quản trị viên.
 \item Lưu công việc, cấu hình, phản hồi và thống kê trong MongoDB.
 \item Tích hợp AI để phân tích, chia nhỏ và lập lại lịch công việc.
 \item Nhắc hạn và đọc nội dung bằng vietTTS hoặc giọng đọc trên thiết bị.
 \item Cung cấp bảng điều khiển quản trị và luồng tiếp nhận phản hồi.
\end{itemize}

\section{Phạm vi}
Đề tài tập trung vào quản lý công việc cá nhân. Hệ thống chưa triển khai cộng tác nhóm,
biểu đồ Gantt, đồng bộ lịch ngoài hoặc nhận dạng giọng nói. Voice hiện tại là chuyển văn
bản thành tiếng nói (TTS), không phải chuyển giọng nói thành văn bản (STT).

\section{Phương pháp thực hiện}
Nhóm áp dụng khảo sát yêu cầu, mô hình hóa kiến trúc, phát triển lặp, kiểm thử tự động
và kiểm chứng trên Android. Nội dung được đối chiếu trực tiếp với mã nguồn tại
\url{https://github.com/Hau-altf4/action_planner}.

\chapter{CƠ SỞ LÝ THUYẾT VÀ CÔNG NGHỆ}
\section{Quản lý công việc và Time-Blocking}
Time-Blocking chia ngày thành các khoảng dành riêng cho từng nhiệm vụ. Trong hệ thống,
mỗi công việc có thời lượng ước tính; lịch bắt đầu được xác định và các công việc liên
tiếp cách nhau 10 phút. Người dùng có thể yêu cầu AI đề xuất lại thứ tự, xem trước và
chỉ lưu phương án sau khi xác nhận.

\section{Flutter và quản lý trạng thái}
Flutter cung cấp widget đa nền tảng. Dự án sử dụng \texttt{AppState} kế thừa
\texttt{ChangeNotifier}; các màn hình nhận cùng thể hiện trạng thái qua hàm khởi tạo.
Cách tổ chức này phù hợp quy mô đồ án và không phụ thuộc gói Provider bên ngoài.

\section{REST, JWT và MongoDB}
Flutter giao tiếp với Express qua JSON/HTTP. JWT chứa định danh, email và vai trò;
middleware xác minh token trước khi truy cập dữ liệu. Mật khẩu được băm bằng bcrypt.
Dự án dùng MongoDB Native Driver trực tiếp, không dùng Mongoose.

\section{AI đa nhà cung cấp}
Máy chủ chuẩn hóa bốn nhà cung cấp: Ollama, OpenAI, Gemini và Anthropic. Ollama là
lựa chọn mặc định chạy cục bộ; cấu hình quản trị hoặc cấu hình riêng của người dùng có
thể ghi đè. Kết quả AI được làm sạch và kiểm tra trước khi chuyển thành dữ liệu ứng dụng.

\section{Text-to-Speech tiếng Việt}
vietTTS chuyển văn bản tiếng Việt thành dạng sóng âm. Backend gọi Python bằng danh
sách tham số an toàn, sinh WAV trong thư mục tạm và xóa sau khi trả kết quả. Khi Python
hoặc mô hình chưa được cài, Flutter chuyển sang \texttt{flutter\_tts} trên thiết bị.

\chapter{PHÂN TÍCH YÊU CẦU}
\section{Tác nhân}
\begin{table}[H]\centering\caption{Tác nhân của hệ thống}
\begin{tabularx}{\textwidth}{p{3.2cm}X}\toprule
\textbf{Tác nhân} & \textbf{Quyền và trách nhiệm}\\\midrule
Người dùng & Đăng ký, đăng nhập, quản lý công việc, lịch, AI, voice, phản hồi và cấu hình AI cá nhân.\\
Quản trị viên & Xem chỉ số, người dùng, phản hồi, trả lời và cấu hình AI mặc định.\\
Dịch vụ AI & Phân tích công việc, hội thoại và đề xuất lịch.\\
Bộ máy TTS & Chuyển nội dung thành WAV hoặc phát bằng giọng thiết bị.\\\bottomrule
\end{tabularx}\end{table}

\section{Yêu cầu chức năng}
\begin{longtable}{p{1.1cm}p{4.3cm}p{9.4cm}}
\caption{Yêu cầu chức năng đã hiện thực}\label{tab:functions}\\
\toprule\textbf{Mã}&\textbf{Nhóm}&\textbf{Mô tả}\\\midrule\endfirsthead
\toprule\textbf{Mã}&\textbf{Nhóm}&\textbf{Mô tả}\\\midrule\endhead
F01&Tài khoản&Đăng ký, đăng nhập, khôi phục phiên, đăng xuất; vai trò user/admin.\\
F02&Công việc&Tạo, cập nhật, xem, lọc trạng thái và đổi trạng thái.\\
F03&Chi tiết&Danh mục, ưu tiên, thời lượng, hạn, bước nhỏ và gợi ý AI.\\
F04&Lịch&Tạo lịch theo thời lượng, xem và xác nhận phương án lập lại lịch.\\
F05&AI&Hội thoại theo ngữ cảnh, phân tích và chia nhỏ công việc.\\
F06&Thông báo&Nhắc trước hạn 10 phút và cảnh báo tại thời điểm hạn.\\
F07&Giọng nói&Đọc bằng vietTTS; dự phòng bằng TTS của thiết bị.\\
F08&Phản hồi&Gửi đề xuất, khiếu nại, đánh giá và xem trả lời.\\
F09&Quản trị&Chỉ số, người dùng, phản hồi và cấu hình AI mặc định.\\
F10&Cache&Lưu danh sách công việc gần nhất để mở ứng dụng nhanh.\\\bottomrule
\end{longtable}

\section{Yêu cầu phi chức năng}
\begin{itemize}
 \item \textbf{Bảo mật:} bcrypt, JWT hết hạn 7 ngày, AES-256-GCM và không commit \texttt{.env}.
 \item \textbf{Tách dữ liệu:} truy vấn công việc ràng buộc \texttt{userId} từ token.
 \item \textbf{Khả dụng:} cache danh sách và TTS dự phòng trên thiết bị.
 \item \textbf{Bảo trì:} repository tách HTTP khỏi UI; voice backend ở mô-đun riêng.
 \item \textbf{Đa nền tảng:} khung dự án hỗ trợ Android, iOS, web và desktop.
\end{itemize}

\chapter{THIẾT KẾ HỆ THỐNG}
\section{Kiến trúc tổng thể}
\begin{figure}[H]\centering
\begin{tikzpicture}[node distance=1.2cm and 1.1cm,
 box/.style={draw,rounded corners,minimum width=3.2cm,minimum height=1.1cm,align=center,fill=schoolblue!8},
 ext/.style={draw,rounded corners,minimum width=2.9cm,minimum height=.9cm,align=center,fill=schoolred!7},
 flow/.style={-{Latex[length=2.5mm]},thick}]
\node[box](ui){Flutter UI\\Mobile/Web/Desktop};
\node[box,right=of ui](state){AppState + Repository\\HTTP, cache, token};
\node[box,right=of state](api){Node.js + Express\\JWT và nghiệp vụ};
\node[box,below=of api](db){MongoDB\\6 collection};
\node[ext,above left=.8cm and -.2cm of api](ai){Ollama/OpenAI\\Gemini/Anthropic};
\node[ext,above right=.8cm and -.2cm of api](tts){vietTTS\\Python + mô hình};
\draw[flow](ui)--(state);\draw[flow](state)--node[above]{JSON/HTTP}(api);
\draw[flow](api)--(db);\draw[flow](api)--(ai);\draw[flow](api)--(tts);
\draw[flow,dashed,bend left=28](ui)to node[below]{TTS dự phòng}(tts);
\end{tikzpicture}\caption{Kiến trúc triển khai thực tế}\end{figure}

\section{Thiết kế dữ liệu}
\begin{table}[H]\centering\caption{Các collection MongoDB thực tế}
\begin{tabularx}{\textwidth}{p{3.2cm}X}\toprule\textbf{Collection}&\textbf{Nội dung}\\\midrule
\texttt{users}&Hồ sơ, email, passwordHash, role và thời điểm hoạt động.\\
\texttt{tasks}&Công việc, ưu tiên, trạng thái, thời lượng, hạn, lịch, bước và gợi ý AI.\\
\texttt{user\_settings}&Cấu hình AI riêng và khóa đã mã hóa.\\
\texttt{app\_settings}&Cấu hình AI mặc định của quản trị viên.\\
\texttt{feedback}&Loại, nội dung, đánh giá, trạng thái và câu trả lời.\\
\texttt{ai\_usage}&Nhà cung cấp, mô hình, độ trễ; tự hết hạn sau 90 ngày.\\\bottomrule
\end{tabularx}\end{table}

\begin{table}[H]\centering\caption{Các trường chính của document tasks}
\begin{tabularx}{\textwidth}{p{3.1cm}p{3cm}X}\toprule\textbf{Trường}&\textbf{Kiểu}&\textbf{Ý nghĩa}\\\midrule
\texttt{userId}&ObjectId&Chủ sở hữu lấy từ JWT.\\
\texttt{title}, \texttt{description}&String&Tiêu đề và mô tả.\\
\texttt{category}&String&Danh mục hiển thị.\\
\texttt{priority}&String&low, medium, high, urgent.\\
\texttt{status}&String&pending, inProgress, completed.\\
\texttt{estimatedMinutes}&Number&Thời lượng dự kiến.\\
\texttt{deadline}&Date&Hạn hoàn thành.\\
\texttt{scheduledStart}&Date/null&Thời điểm bắt đầu.\\
\texttt{steps}&Array&Các bước con gồm id, title, status.\\
\texttt{aiSuggestion}&String/null&Gợi ý AI.\\\bottomrule
\end{tabularx}\end{table}

\section{Thiết kế API}
\small
\begin{longtable}{p{1.5cm}p{5.2cm}p{7.4cm}}
\caption{Các endpoint được triển khai}\label{tab:api}\\
\toprule\textbf{HTTP}&\textbf{Endpoint}&\textbf{Chức năng}\\\midrule\endfirsthead
\toprule\textbf{HTTP}&\textbf{Endpoint}&\textbf{Chức năng}\\\midrule\endhead
GET&\texttt{/health}&Kiểm tra MongoDB, AI và voice.\\
POST&\texttt{/api/auth/register}&Đăng ký và cấp JWT.\\
POST&\texttt{/api/auth/login}&Đăng nhập và cấp JWT.\\
GET&\texttt{/api/auth/me}&Khôi phục hồ sơ.\\
GET/POST&\texttt{/api/tasks}&Lấy danh sách/tạo công việc.\\
PUT&\texttt{/api/tasks/:id}&Cập nhật công việc.\\
PATCH&\texttt{/api/tasks/:id/status}&Cập nhật trạng thái.\\
POST&\texttt{/api/ai/chat}&Hội thoại AI.\\
POST&\texttt{/api/ai/analyze-task}&Phân tích công việc.\\
POST&\texttt{/api/plans/replan}&Tạo phương án lịch.\\
POST&\texttt{/api/plans/replan/accept}&Xác nhận phương án.\\
GET/PUT&\texttt{/api/settings/ai}&Đọc/lưu AI riêng.\\
POST&\texttt{/api/settings/ai/test}&Kiểm tra AI.\\
POST/GET&\texttt{/api/feedback}&Gửi/xem phản hồi.\\
POST&\texttt{/api/voice/synthesize}&Tạo WAV bằng vietTTS.\\
GET&\texttt{/api/admin/dashboard}&Chỉ số quản trị.\\
GET&\texttt{/api/admin/users}&Danh sách người dùng.\\
GET/PATCH&\texttt{/api/admin/feedback[/:id]}&Quản lý phản hồi.\\
GET/PUT&\texttt{/api/admin/settings/ai}&AI mặc định.\\\bottomrule
\end{longtable}\normalsize

\chapter{CÀI ĐẶT VÀ HIỆN THỰC}
\section{Cấu trúc mã nguồn}
\begin{lstlisting}[caption={Cấu trúc rút gọn của repository}]
action_planner/
|-- lib/
|   |-- data/task_repository.dart
|   |-- models/       |-- services/
|   |-- state/app_state.dart
|   `-- ui/
|-- backend/src/server.js
|-- backend/src/voice.js
|-- android/ ios/ web/ windows/ linux/ macos/
|-- test/widget_test.dart
`-- report/overleaf/
\end{lstlisting}

\section{Đăng nhập và dữ liệu}
Ứng dụng thử các địa chỉ backend phù hợp thiết bị, gọi endpoint đăng nhập, lưu JWT
trong \texttt{flutter\_secure\_storage} trên native hoặc SharedPreferences trên web.
Khi mở lại, ứng dụng gọi \texttt{/api/auth/me}; token lỗi được xóa. Repository ánh xạ
JSON thành \texttt{PlannerTask}. Danh sách máy chủ trả về được cache cục bộ.

\section{Lịch và thông báo}
Lịch chỉ chứa công việc chưa hoàn thành và ưu tiên \texttt{scheduledStart} đã lưu.
Ứng dụng dùng \texttt{flutter\_local\_notifications} và timezone, lên lịch nhắc trước
hạn 10 phút và cảnh báo tại thời điểm hạn. Khi hoàn thành, thông báo được hủy/lập lại.

\section{Bảo mật khóa AI}
Khóa API không được trả lại client. Backend băm bí mật cấu hình thành khóa 256 bit,
mã hóa AES-256-GCM với IV ngẫu nhiên 12 byte, lưu IV, authentication tag và ciphertext.
Màn hình chỉ nhận cờ \texttt{hasApiKey}.

\placeholderfigure{images/dang-nhap.png}{Màn hình đăng nhập}
\placeholderfigure{images/danh-sach-cong-viec.png}{Danh sách công việc và nút voice}
\placeholderfigure{images/chinh-sua-cong-viec.png}{Màn hình tạo và phân tích công việc}

\chapter{TRỢ LÝ AI VÀ GIỌNG NÓI}
\section{Luồng xử lý AI}
\begin{figure}[H]\centering
\begin{tikzpicture}[node distance=.8cm,
 flowstep/.style={draw,rounded corners,fill=schoolblue!7,text width=11.8cm,minimum height=.85cm,align=center},
 flow/.style={-{Latex[length=2.5mm]},thick}]
\node[flowstep](a){Người dùng nhập câu hỏi hoặc yêu cầu phân tích};
\node[flowstep,below=of a](b){Flutter gửi JWT, lịch sử và dữ liệu cần thiết};
\node[flowstep,below=of b](c){Backend chọn cấu hình cá nhân, quản trị hoặc môi trường};
\node[flowstep,below=of c](d){Gọi Ollama/OpenAI/Gemini/Anthropic và ghi nhận độ trễ};
\node[flowstep,below=of d](e){Chuẩn hóa kết quả và trả về Flutter};
\draw[flow](a)--(b);\draw[flow](b)--(c);\draw[flow](c)--(d);\draw[flow](d)--(e);
\end{tikzpicture}\caption{Pipeline trợ lý AI đa nhà cung cấp}\end{figure}

Endpoint phân tích trả mức ưu tiên, thời lượng, bước và gợi ý. Lập lại lịch lấy công
việc chưa hoàn thành, tạo phương án xem trước và chỉ lưu khi người dùng chọn áp dụng.
Nếu AI không khả dụng, backend có logic dự phòng dựa trên ưu tiên và thời hạn.

\section{Thiết kế voice}
\begin{figure}[H]\centering
\begin{tikzpicture}[node distance=1cm and 1cm,
 box/.style={draw,rounded corners,fill=schoolred!6,minimum width=3cm,minimum height=.9cm,align=center},
 flow/.style={-{Latex[length=2.5mm]},thick}]
\node[box](tap){Bấm nút loa};\node[box,right=of tap](post){POST voice\\JWT + văn bản};
\node[box,right=of post](engine){vietTTS\\tạo WAV};
\node[box,below=of engine](fallback){TTS hệ thống\\vi-VN};
\node[box,left=of fallback](play){Phát và dừng\\âm thanh};
\draw[flow](tap)--(post);\draw[flow](post)--(engine);\draw[flow](engine)--(play);
\draw[flow,dashed](post)--node[right]{lỗi/503}(fallback);\draw[flow](fallback)--(play);
\end{tikzpicture}\caption{Luồng vietTTS và dự phòng trên thiết bị}\end{figure}

Nội dung giới hạn 500 ký tự. Backend dùng \texttt{spawn} với \texttt{shell:false},
giới hạn thời gian, đọc WAV rồi xóa thư mục tạm. Flutter phát WAV bằng
\texttt{audioplayers}. Nếu backend lỗi, \texttt{flutter\_tts} đọc với locale
\texttt{vi-VN}. vietTTS chỉ là TTS; tạo công việc bằng giọng nói cần thêm STT và bước
xác nhận, được xác định là hướng phát triển chứ chưa phải chức năng hoàn thành.

\chapter{KIỂM THỬ VÀ ĐÁNH GIÁ}
\section{Kết quả có thể tái lập}
\begin{table}[H]\centering\caption{Kết quả kiểm tra}
\begin{tabularx}{\textwidth}{p{4cm}p{3cm}X}\toprule
\textbf{Kiểm tra}&\textbf{Kết quả}&\textbf{Lệnh/minh chứng}\\\midrule
Phân tích tĩnh Flutter&Đạt, 0 lỗi&\texttt{flutter analyze}\\
Widget/unit test&Đạt, 7/7&\texttt{flutter test}\\
Cú pháp backend&Đạt&\texttt{node --check backend/src/server.js}\\
Cú pháp voice&Đạt&\texttt{node --check backend/src/voice.js}\\
Biên dịch Android&Đạt&\texttt{flutter build apk --debug}\\
Health backend&Đạt&MongoDB connected, AI configured; voice phản ánh đúng trạng thái.\\\bottomrule
\end{tabularx}\end{table}

\section{Bảy kịch bản tự động}
\begin{enumerate}
 \item Lịch cố định và giữ nguyên sau vòng chuyển JSON.
 \item Hiển thị màn hình đăng nhập đúng cấu trúc.
 \item Nhập tiếng Việt và không tự điền tài khoản.
 \item Câu trả lời AI cập nhật giao diện an toàn.
 \item Đăng nhập không áp dụng nhầm quy tắc đăng ký.
 \item Rời trình chỉnh sửa khi AI chạy không gọi setState sau dispose.
 \item Yêu cầu sửa phương án lịch thay thế hộp thoại an toàn.
\end{enumerate}

\section{Đánh giá trung thực}
Hệ thống đáp ứng các luồng chính và có dự phòng cho AI/voice. Báo cáo không khẳng
định 120 ms, 60 fps, 80 MB RAM hay độ chính xác AI nếu chưa có benchmark và dữ liệu
đo đi kèm. Kết quả chỉ sử dụng minh chứng tái lập được.

\chapter{KẾT LUẬN VÀ HƯỚNG PHÁT TRIỂN}
\section{Kết luận}
Đồ án xây dựng sản phẩm xuyên suốt từ Flutter đến REST API và MongoDB, có xác thực,
phân quyền, công việc, lịch, thông báo, AI đa nhà cung cấp, quản trị, phản hồi và TTS
tiếng Việt. Kiến trúc repository và cấu hình môi trường giúp dễ thay AI hoặc bộ máy TTS.

\section{Hạn chế}
\begin{itemize}
 \item Chưa có xóa công việc, tìm kiếm toàn văn và cộng tác nhóm.
 \item Cache phục vụ đọc nhanh; chưa có hàng đợi đồng bộ ngoại tuyến hai chiều.
 \item vietTTS cần Python/mô hình; giọng dự phòng phụ thuộc thiết bị.
 \item Chưa có benchmark tự động cho API, FPS và RAM.
 \item Chưa có kiểm thử tích hợp backend với MongoDB tạm trong CI.
\end{itemize}

\section{Hướng phát triển}
\begin{itemize}
 \item Thêm xóa mềm, tìm kiếm, phân trang và xử lý xung đột đồng bộ.
 \item Bổ sung STT để tạo công việc bằng giọng nói và xác nhận trước khi lưu.
 \item Đồng bộ Google/Apple Calendar và hỗ trợ công việc lặp lại.
 \item Tạo CI chạy phân tích, test Flutter và test API trên mỗi commit.
 \item Xây dựng benchmark tái lập và theo dõi độ trễ AI.
 \item Đóng gói vietTTS thành dịch vụ hoặc chuyển sang LightSpeed khi phù hợp.
\end{itemize}

\appendix
\chapter{CÀI ĐẶT VÀ TÁI LẬP}
\section{Backend}
\begin{lstlisting}[language=bash,caption={Khởi tạo backend}]
cd backend
copy .env.example .env
npm install
npm run dev
\end{lstlisting}
Các biến tối thiểu gồm \texttt{JWT\_SECRET}, \texttt{MONGODB\_URI} và
\texttt{MONGODB\_DB}. Không commit \texttt{backend/.env}.

\section{Flutter}
\begin{lstlisting}[language=bash,caption={Chạy và kiểm tra Flutter}]
flutter pub get
flutter analyze
flutter test
flutter run --dart-define=API_BASE_URL=http://127.0.0.1:3000
\end{lstlisting}
Android Emulator dùng \texttt{10.0.2.2}; điện thoại thật có thể dùng
\texttt{adb reverse tcp:3000 tcp:3000} hoặc địa chỉ LAN.

\section{vietTTS}
Khai báo \texttt{VIETTTS\_DIR}, \texttt{VIETTTS\_PYTHON},
\texttt{VIETTTS\_LEXICON} và \texttt{VIETTTS\_TIMEOUT\_MS}. Khi chưa cấu hình,
health trả \texttt{voice: not\_configured} và Flutter dùng TTS thiết bị.

\chapter{DÙNG TRÊN OVERLEAF}
\begin{enumerate}
 \item Tải toàn bộ thư mục \texttt{report/overleaf} lên Overleaf.
 \item Chọn \textbf{XeLaTeX} tại \textbf{Menu -- Compiler}.
 \item Sửa \texttt{\string\studentTwoId} ở đầu \texttt{main.tex}.
 \item Thêm logo/ảnh vào \texttt{images} theo đúng tên gợi ý nếu có.
 \item Chọn \texttt{main.tex} làm Main document và Recompile.
\end{enumerate}

\begin{thebibliography}{99}
\bibitem{flutter} Google, \textit{Flutter Documentation}. \url{https://docs.flutter.dev/}
\bibitem{dart} Google, \textit{Dart Documentation}. \url{https://dart.dev/guides}
\bibitem{node} OpenJS Foundation, \textit{Node.js Documentation}. \url{https://nodejs.org/docs/latest/api/}
\bibitem{express} Express.js, \textit{Express API Reference}. \url{https://expressjs.com/}
\bibitem{mongodb} MongoDB Inc., \textit{MongoDB Manual}. \url{https://www.mongodb.com/docs/manual/}
\bibitem{jwt} IETF, \textit{RFC 7519: JSON Web Token}. \url{https://www.rfc-editor.org/rfc/rfc7519}
\bibitem{ollama} Ollama, \textit{Ollama API}. \url{https://github.com/ollama/ollama/blob/main/docs/api.md}
\bibitem{openai} OpenAI, \textit{Responses API}. \url{https://platform.openai.com/docs/api-reference/responses}
\bibitem{viettts} NTT123, \textit{vietTTS}. \url{https://github.com/NTT123/vietTTS}
\bibitem{fluttertts} flutter\_tts contributors, \textit{Flutter TTS Plugin}. \url{https://pub.dev/packages/flutter_tts}
\end{thebibliography}
\end{document}
